\section{\texttt{piret.py}: function-by-function implementation}
Read public differential helpers, external target adapters, SM4 inversion/filtering/peeling, AES quartet filtering and reverse schedules before terminal formatting. Keep the legacy known-key relation lesson separate from the ciphertext-only recovery paths.

\subsection{Implementation topics and presentation order}
\begin{enumerate}
\item \textbf{Experimental roles.} Separate the normal cipher, key-owning simulated target, ciphertext-based solver and teaching display. Explain why showing a hidden state for learning does not permit giving it to the solver.
\item \textbf{Differential helpers.} Explain \texttt{difference}, \texttt{clone}, \texttt{gf\_multiply} and \texttt{predicted\_mixcolumns\_difference}. For a row-zero fault, derive the four-byte pattern $(2e,e,e,3e)$ in the AES field.
\item \textbf{AES fault injection.} Follow \texttt{encrypt\_with\_trace} and its observer before MixColumns. Compare round nine, which reaches one quartet, with round eight, which reaches all four groups. \texttt{trace\_round8} is a compatibility name that also accepts an explicit round-nine request.
\item \textbf{AES candidate construction.} Explain \texttt{output\_positions} and \texttt{key\_candidates}: undo the final key XOR/S-box under each guess, match a shared error row and magnitude, and retain correlated four-byte tuples. Independent byte guesses are insufficient.
\item \textbf{AES complete key recovery.} Explain per-column intersections in the round-nine lesson and \texttt{recover\_last\_key} for the round-eight path. Once K10 is unique, follow \texttt{reverse\_key\_schedule} to K0 and verify by re-encryption. Keep the known-key \texttt{relation\_demo} distinct from unknown-key reconstruction.
\item \textbf{SM4 controlled-byte model.} Follow \texttt{sm4\_fault\_encrypt}, \texttt{inverse\_rows}, \texttt{inverse\_linear} and \texttt{sm4\_byte\_candidates}. Undo L, test S-box input candidates using the known byte mask, and convert them to subkey-byte guesses using the observable words.
\item \textbf{SM4 multi-round recovery.} Follow \texttt{peel\_sm4} and \texttt{recover\_sm4\_master\_key}. Recover $rk_{31}$ through $rk_{28}$ using fresh faults at each target round and only previously recovered later keys. Explain the additional selectable-round fault assumption.
\item \textbf{Master-key reconstruction and output.} Follow \texttt{reverse\_sm4}, remove FK and assemble the original master key. Earlier SM4 subkeys are derived by the reversed schedule, whereas the final four came from DFA. Explain \texttt{Terminal}, \texttt{main}, red differences, step mode and shared TXT/HTML logging.
\item \textbf{Ambiguity and failure.} Distinguish an incompatible empty intersection from multiple compatible survivors. Neither case permits filling missing bytes from the teaching key or claiming successful complete recovery.
\end{enumerate}
\noindent\textbf{Suggested explanation sequence:} fault location $\rightarrow$ observed difference $\rightarrow$ candidate equations $\rightarrow$ intersections $\rightarrow$ recovered subkeys $\rightarrow$ inverse schedule $\rightarrow$ verified master key.

\subsection{\texttt{difference}}
\begin{lightcode}
def difference(a, b):
\end{lightcode}
\noindent\textbf{Source:} complete \texttt{piret.py} listing in Part 8, lines 65--67.
\begin{enumerate}
\item Accept two AES matrices with matching coordinates.
\item Return a new matrix containing the coordinate-wise XOR. The inputs remain unchanged, allowing a teacher to compare saved states safely.
\end{enumerate}
\subsection{\texttt{clone}}
\begin{lightcode}
def clone(state):
\end{lightcode}
\noindent\textbf{Source:} complete \texttt{piret.py} listing in Part 8, lines 70--72.
\begin{enumerate}
\item Accept a nested row-list state.
\item Copy every row with slicing and return the new outer list. Copying only the outer list would leave shared rows and would not protect a saved snapshot against later in-place changes.
\end{enumerate}
\subsection{\texttt{gf\_\allowbreak{}multiply}}
\begin{lightcode}
def gf_multiply(a, b):
\end{lightcode}
\noindent\textbf{Source:} complete \texttt{piret.py} listing in Part 8, lines 75--81.
\begin{enumerate}
\item Accept two field bytes under the historical attack-helper name.
\item Return aes.gf\_mul(a,b). Sharing the cipher's field arithmetic avoids a second implementation with a different reduction polynomial.
\end{enumerate}
\subsection{\texttt{predicted\_\allowbreak{}mixcolumns\_\allowbreak{}difference}}
\begin{lightcode}
def predicted_mixcolumns_difference(row, fault):
\end{lightcode}
\noindent\textbf{Source:} complete \texttt{piret.py} listing in Part 8, lines 84--89.
\begin{enumerate}
\item Accept the faulty input row and a nonzero byte difference.
\item Multiply that scalar error by the matching column of MIX\_MATRIX using gf\_multiply. Return four predicted differences; row zero gives the ordered pattern [2e,e,e,3e].
\end{enumerate}
\subsection{\texttt{reverse\_\allowbreak{}last\_\allowbreak{}round\_\allowbreak{}difference}}
\begin{lightcode}
def reverse_last_round_difference(correct_ciphertext, faulty_ciphertext, last_round_key,):
\end{lightcode}
\noindent\textbf{Source:} complete \texttt{piret.py} listing in Part 8, lines 92--118.
\begin{enumerate}
\item Accept correct/faulty ciphertexts and an already known K10 matrix, then convert the two blocks to state matrices. Remove final AddRoundKey, undo ShiftRows and apply INV\_SBOX.
\item Return the XOR of the resulting states. This verifies a known-key differential relation; because K10 is supplied, it is not the unknown-key candidate solver.
\end{enumerate}
\subsection{\texttt{encrypt\_\allowbreak{}with\_\allowbreak{}trace}}
\begin{lightcode}
def encrypt_with_trace(plaintext, key, row=0, column=0, mask=0, fault_round=8):
\end{lightcode}
\noindent\textbf{Source:} complete \texttt{piret.py} listing in Part 8, lines 126--144.
\begin{enumerate}
\item Validate coordinates, byte mask and the choice of round eight or nine. Install a local observer that changes one state byte before the selected MixColumns event.
\item Run aes.aes\_encrypt once with this observer and return its ciphertext plus immutable after-operation snapshots. The target adapter sees the simulation key, but the candidate solver is called separately with output observations.
\end{enumerate}
\subsection{\texttt{encrypt\_\allowbreak{}with\_\allowbreak{}trace.observer}}
\begin{lightcode}
def observer(number, operation, state):
\end{lightcode}
\noindent\textbf{Source:} complete \texttt{piret.py} listing in Part 8, lines 136--141.
\begin{enumerate}
\item Match the requested round and Before MixColumns event; XOR the selected live state byte. Save a serialized Fault point and after-operation snapshots. Return values are unused because AES observers mutate the matrix.
\end{enumerate}
\subsection{\texttt{sm4\_\allowbreak{}fault\_\allowbreak{}encrypt}}
\begin{lightcode}
def sm4_fault_encrypt(plaintext, key, byte=None, mask=0, fault_round=31):
\end{lightcode}
\noindent\textbf{Source:} complete \texttt{piret.py} listing in Part 8, lines 147--169.
\begin{enumerate}
\item Validate the selected SM4 byte, mask and round; expand the target key only for the diagnostic trace. Install a callback at the Sbox input event that shifts the byte mask into the chosen big-endian position.
\item Run sm4.sm4\_encrypt\_block with this observer and return ciphertext and per-round A, B, linear, new-word and subkey trace fields. Full-recovery oracle closures discard the trace and expose only the ciphertext to the solver.
\end{enumerate}
\subsection{\texttt{sm4\_\allowbreak{}fault\_\allowbreak{}encrypt.observer}}
\begin{lightcode}
def observer(number, operation, value):
\end{lightcode}
\noindent\textbf{Source:} complete \texttt{piret.py} listing in Part 8, lines 161--166.
\begin{enumerate}
\item At the selected Sbox input, XOR the shifted mask into the integer value. Store selected scalar event values in the diagnostic round dictionary and return the possibly replaced value to the SM4 engine.
\end{enumerate}
\subsection{\texttt{trace\_\allowbreak{}round8}}
\begin{lightcode}
def trace_round8(plaintext, key, row=0, column=0, mask=0, fault_round=8):
\end{lightcode}
\noindent\textbf{Source:} complete \texttt{piret.py} listing in Part 8, lines 172--178.
\begin{enumerate}
\item Accept the same explicit parameters as encrypt\_with\_trace, including fault\_round.
\item Delegate directly and return its ciphertext/trace pair. The historical function name does not force round eight: callers can explicitly request round nine.
\end{enumerate}
\subsection{\texttt{encrypt}}
\begin{lightcode}
def encrypt(plaintext, key, inject_fault=False):
\end{lightcode}
\noindent\textbf{Source:} complete \texttt{piret.py} listing in Part 8, lines 181--192.
\begin{enumerate}
\item Read the legacy fault-coordinate globals and choose their mask only when inject\_fault is true. Run the common AES target adapter at round nine.
\item Translate selected trace events into named matrix snapshots and return ciphertext, snapshots and expanded key matrices. This compatibility interface supports the known-key relation lesson, not a ciphertext-only recovery claim.
\end{enumerate}
\subsection{\texttt{aes\_\allowbreak{}target\_\allowbreak{}encrypt}}
\begin{lightcode}
def aes_target_encrypt(plaintext, key):
\end{lightcode}
\noindent\textbf{Source:} complete \texttt{piret.py} listing in Part 8, lines 195--197.
\begin{enumerate}
\item Accept the plaintext and target key.
\item Call aes.aes\_encrypt without an observer and return the clean ciphertext. This supplies an ordinary baseline without changing global cipher settings.
\end{enumerate}
\subsection{\texttt{inverse\_\allowbreak{}rows}}
\begin{lightcode}
def inverse_rows():
\end{lightcode}
\noindent\textbf{Source:} complete \texttt{piret.py} listing in Part 8, lines 208--218.
\begin{enumerate}
\item Construct the binary matrix of SM4 L by applying linear to each of the thirty-two unit-bit words. Augment its rows with an identity matrix.
\item Perform binary Gaussian elimination using row swaps and XOR, then return the rows of the inverse matrix. The lru\_cache decorator stores this fixed result so later calls do not repeat elimination.
\end{enumerate}
\subsection{\texttt{inverse\_\allowbreak{}linear}}
\begin{lightcode}
def inverse_linear(value):
\end{lightcode}
\noindent\textbf{Source:} complete \texttt{piret.py} listing in Part 8, lines 221--222.
\begin{enumerate}
\item Accept a 32-bit output difference.
\item For each cached inverse-matrix row, compute the parity of row AND value and place that parity into its output bit. Return L-inverse(value), using binary arithmetic rather than a guessed inverse rotation.
\end{enumerate}
\subsection{\texttt{sm4\_\allowbreak{}byte\_\allowbreak{}candidates}}
\begin{lightcode}
def sm4_byte_candidates(correct, faulty, byte, mask):
\end{lightcode}
\noindent\textbf{Source:} complete \texttt{piret.py} listing in Part 8, lines 225--239.
\begin{enumerate}
\item Validate the two output blocks, byte index and known mask; require the other three words to agree. Invert the first-word XOR through inverse\_linear and require support in exactly the requested S-box byte.
\item Enumerate all input bytes a satisfying the S-box differential, convert them to key-byte candidates by XOR with the observable three-word combination, and return candidates, input guesses and the 32-bit inverse-L difference.
\end{enumerate}
\subsection{\texttt{reverse\_\allowbreak{}sm4}}
\begin{lightcode}
def reverse_sm4(final_keys, observer=None):
\end{lightcode}
\noindent\textbf{Source:} complete \texttt{piret.py} listing in Part 8, lines 242--252.
\begin{enumerate}
\item Require four consecutive recovered subkeys in rk28,rk29,rk30,rk31 order and place them at schedule indices thirty-two through thirty-five.
\item Iterate backward using CK, sbox and linear\_key to derive earlier schedule words. Undo the four FK constants and pack the resulting master words; optionally report each recovered schedule word to an observer.
\end{enumerate}
\subsection{\texttt{peel\_\allowbreak{}sm4}}
\begin{lightcode}
def peel_sm4(ciphertext, target_round, recovered_keys):
\end{lightcode}
\noindent\textbf{Source:} complete \texttt{piret.py} listing in Part 8, lines 255--271.
\begin{enumerate}
\item Validate the requested round in twenty-eight through thirty-one and reverse ciphertext word order into a rolling window.
\item Undo every later round using only entries already present in recovered\_keys; reject a missing key. Return a reversed-word view at the target round, allowing reuse of the same last-round candidate equations.
\end{enumerate}
\subsection{\texttt{recover\_\allowbreak{}sm4\_\allowbreak{}master\_\allowbreak{}key}}
\begin{lightcode}
def recover_sm4_master_key(correct, fault_oracle, masks=(1, 2, 3, 4, 5, 6, 7, 8), start_byte=0, candidate_filter=None, observer=None):
\end{lightcode}
\noindent\textbf{Source:} complete \texttt{piret.py} listing in Part 8, lines 274--318.
\begin{enumerate}
\item Validate distinct nonzero masks and iterate target rounds thirty-one down to twenty-eight. For each target byte, query the output-only oracle, peel correct and faulty outputs and apply the known-byte filter or the supplied replacement filter.
\item Intersect candidates, emit history events and stop with no master key when ambiguity remains. Once all four final subkeys are uniquely recovered, call reverse\_sm4 and return the master key, recovered subkeys, counts and history.
\end{enumerate}
\subsection{\texttt{output\_\allowbreak{}positions}}
\begin{lightcode}
def output_positions(column):
\end{lightcode}
\noindent\textbf{Source:} complete \texttt{piret.py} listing in Part 8, lines 322--327.
\begin{enumerate}
\item Validate an AES column index in zero through three.
\item Return indices 4*((column-row) modulo 4)+row for rows zero through three. This undoes the final ShiftRows permutation for a quartet while retaining the original row order used by the mixing equation.
\end{enumerate}
\subsection{\texttt{key\_\allowbreak{}candidates}}
\begin{lightcode}
def key_candidates(correct, faulty, column):
\end{lightcode}
\noindent\textbf{Source:} complete \texttt{piret.py} listing in Part 8, lines 330--355.
\begin{enumerate}
\item For each of the four mapped ciphertext bytes, test all 256 K10 byte guesses and group them by the inverse-S-box difference they produce.
\item Enumerate one shared faulty row and nonzero error, use the corresponding four MixColumns coefficients to select compatible lists, and take their Cartesian product. Return a set of correlated four-byte tuples rather than independent byte sets.
\end{enumerate}
\subsection{\texttt{recover\_\allowbreak{}last\_\allowbreak{}key}}
\begin{lightcode}
def recover_last_key(correct, faulty_outputs, observer=None):
\end{lightcode}
\noindent\textbf{Source:} complete \texttt{piret.py} listing in Part 8, lines 358--386.
\begin{enumerate}
\item Process round-eight faulty outputs, requiring every ciphertext byte to change under that specific model. Use key\_candidates for each of four output groups and intersect its tuples across observations.
\item Reject empty intersections; when every group is unique, map all tuple bytes into a sixteen-byte K10 and return it with the observation count. If observations end first, return None instead of a guessed key.
\end{enumerate}
\subsection{\texttt{reverse\_\allowbreak{}key\_\allowbreak{}schedule}}
\begin{lightcode}
def reverse_key_schedule(last_key, observer=None):
\end{lightcode}
\noindent\textbf{Source:} complete \texttt{piret.py} listing in Part 8, lines 394--415.
\begin{enumerate}
\item Place the recovered sixteen K10 bytes into words forty through forty-three. For descending indices, copy the previous word, applying RotWord, forward SubWord and Rcon only when the index is a multiple of four.
\item XOR this transformed word with the current word to obtain the word four positions earlier. Optionally notify a presentation observer, then concatenate the first four words to return K0.
\end{enumerate}
\subsection{\texttt{show\_\allowbreak{}state}}
\begin{lightcode}
def show_state(title, state):
\end{lightcode}
\noindent\textbf{Source:} complete \texttt{piret.py} listing in Part 8, lines 423--428.
\begin{enumerate}
\item Accept a title and a row-indexed AES matrix.
\item Print a heading and format each matrix row as two-digit hexadecimal bytes. This is presentation only and returns no cryptographic result.
\end{enumerate}
\subsection{\texttt{show\_\allowbreak{}difference}}
\begin{lightcode}
def show_difference(title, a, b):
\end{lightcode}
\noindent\textbf{Source:} complete \texttt{piret.py} listing in Part 8, lines 431--433.
\begin{enumerate}
\item Accept a title and two saved matrices.
\item Compute their difference and pass it to show\_state. Neither input is changed, and no candidate filtering occurs.
\end{enumerate}
\subsection{\texttt{Terminal.\_\allowbreak{}\_\allowbreak{}init\_\allowbreak{}\_\allowbreak{}}}
\begin{lightcode}
def __init__(self, color="auto", step=False):
\end{lightcode}
\noindent\textbf{Source:} complete \texttt{piret.py} listing in Part 8, lines 438--447.
\begin{enumerate}
\item Resolve color mode from the explicit option, terminal detection and NO\_COLOR environment variable, and store the step-mode flag.
\item When color is active on Windows, enable virtual-terminal processing through the console API. Presentation configuration does not change target bytes or solver equations.
\end{enumerate}
\subsection{\texttt{Terminal.paint}}
\begin{lightcode}
def paint(self, text, code="36"):
\end{lightcode}
\noindent\textbf{Source:} complete \texttt{piret.py} listing in Part 8, lines 449--451.
\begin{enumerate}
\item Accept text and an ANSI style code.
\item Return the styled string plus reset sequence when color is enabled, or plain text otherwise. Byte values are formatted elsewhere before this method receives them.
\end{enumerate}
\subsection{\texttt{Terminal.section}}
\begin{lightcode}
def section(self, title):
\end{lightcode}
\noindent\textbf{Source:} complete \texttt{piret.py} listing in Part 8, lines 453--456.
\begin{enumerate}
\item Accept a teaching-section title.
\item Print a separator and the title through paint. This creates a visual boundary without advancing an algorithm or changing candidate sets.
\end{enumerate}
\subsection{\texttt{Terminal.pause}}
\begin{lightcode}
def pause(self):
\end{lightcode}
\noindent\textbf{Source:} complete \texttt{piret.py} listing in Part 8, lines 458--461.
\begin{enumerate}
\item Inspect the stored step-mode flag.
\item Call input only when stepping is enabled; otherwise return immediately. Batch and default execution therefore remain noninteractive.
\end{enumerate}
\subsection{\texttt{Terminal.compare}}
\begin{lightcode}
def compare(self, correct, faulty):
\end{lightcode}
\noindent\textbf{Source:} complete \texttt{piret.py} listing in Part 8, lines 463--477.
\begin{enumerate}
\item Accept two serialized states, compute their XOR, and display the bytes in AES row order using index 4*column+row.
\item Color changed faulty bytes and nonzero differences red, then print changed-byte positions and Hamming distance. These metrics describe the given pair rather than a recovery success probability.
\end{enumerate}
\subsection{\texttt{relation\_\allowbreak{}demo}}
\begin{lightcode}
def relation_demo():
\end{lightcode}
\noindent\textbf{Source:} complete \texttt{piret.py} listing in Part 8, lines 485--585.
\begin{enumerate}
\item Decode legacy constants, validate the fault descriptor and obtain clean and faulted round-nine snapshots with encrypt.
\item Compare observed MixColumns differences with the predicted pattern, then undo the final round using the known K10 and assert the differential relation. Explicitly report that the key was known; this routine does not recover an unknown key.
\end{enumerate}
\subsection{\texttt{main}}
\begin{lightcode}
def main(argv=None):
\end{lightcode}
\noindent\textbf{Source:} complete \texttt{piret.py} listing in Part 8, lines 593--715.
\begin{enumerate}
\item Parse options, load the target configuration and route SM4 requests to its controlled-byte demonstration. For AES, obtain the clean trace and generate reproducible round-eight faulty outputs through a local generator.
\item Call recover\_last\_key using only ciphertexts, validate the recovered differential equations and reverse the schedule. Verify by re-encryption and final simulation ground truth; return zero on success or two on unresolved ambiguity.
\end{enumerate}
\subsection{\texttt{main.observations}}
\begin{lightcode}
def observations():
\end{lightcode}
\noindent\textbf{Source:} complete \texttt{piret.py} listing in Part 8, lines 659--666.
\begin{enumerate}
\item Iterate the reproducible fault-mask sequence, run the target adapter, display selected teaching comparisons and yield only each faulty ciphertext. Lazy yielding lets the solver stop generating faults once all groups are unique.
\end{enumerate}
\subsection{\texttt{main.report}}
\begin{lightcode}
def report(index, column, before, fresh, survivors):
\end{lightcode}
\noindent\textbf{Source:} complete \texttt{piret.py} listing in Part 8, lines 668--675.
\begin{enumerate}
\item Receive the observation index, group, previous count, fresh count and surviving tuples from recover\_last\_key. Print candidate evolution and representative tuples; do not supply any new candidate or key byte.
\end{enumerate}
\subsection{\texttt{main.schedule\_\allowbreak{}report}}
\begin{lightcode}
def schedule_report(index, current, transformed, previous):
\end{lightcode}
\noindent\textbf{Source:} complete \texttt{piret.py} listing in Part 8, lines 700--703.
\begin{enumerate}
\item Receive the current expanded word, transformed predecessor and recovered earlier word. Format the reverse-schedule XOR equation and optionally pause on a key-schedule core boundary.
\end{enumerate}
